\section{Introduction}
\label{sec:intro}

Hammers --- tools that discharge proof obligations of an
interactive proof assistant by calling automated theorem provers
(ATPs) --- have become standard equipment for proof assistants
based on simple type theory, and increasingly for those based on
dependent type theory: CoqHammer for Coq/Rocq
\cite{CzajkaKaliszyk2018}, and more recently Lean-auto and the Lean
hammer effort \cite{QianEtAl2025,ZhuEtAl2025}.  A hammer has three
phases: premise selection, translation of the goal and premises
into the ATPs' logic, and reconstruction of the found proof inside
the assistant.  This paper is about the middle phase for dependent
type theory, and specifically about two optimisations that are
folklore for simple type theory but have never been given a
rigorous treatment for the Calculus of Inductive Constructions
(CIC):

\begin{itemize}
\item \emph{Heuristic monomorphisation}: replacing premises that
  quantify over types (``$\forall A{:}\mathsf{Type}$.
  $\forall x{:}A$. $\forall l{:}\mathsf{list}\,A$. \ldots'') by
  finitely many instances at the ground types occurring in the problem
  ($A \coloneqq \mathsf{nat}$, \ldots).  For Sledgehammer-style
  hammers this is known to outperform not only the polymorphic
  encodings but even native polymorphic prover support
  \cite{Blanchette2016,Desharnais2022}; for CIC it has remained a
  to-do item.
\item \emph{Guard erasure}: the CoqHammer translation is an untyped
  encoding that retains typing information as \emph{guards} ---
  atoms $\hastype(x, \tau)$ read ``$x$ has type $\tau$'' ---
  attached to every quantifier.  Guards multiply literals and
  clutter every clause.  For many-sorted first-order logic, the
  monotonicity analysis of Claessen, Lillieström and Smallbone
  \cite{Claessen2011} and Blanchette, B\"ohme, Popescu and
  Smallbone \cite{Blanchette2016} identifies which sorts' guards
  can be erased soundly.  Whether and how that analysis transfers
  to dependent type theory was posed as an open question in the
  soundness study of the CoqHammer translation
  \cite[\S6]{Czajka2016}, whose own proof-theoretic method is
  incompatible with the model constructions monotonicity requires.
  Meanwhile the implementation omits certain guards on empirical
  grounds, with a known concrete counterexample delimiting the
  practice \cite[\S5.6]{CzajkaKaliszyk2018}.
\end{itemize}

Both optimisations are \emph{dangerous}: instantiation interacts
with the type-directed translation, and guard erasure strengthens
premises --- erase the guard of a quantifier over a one-element
type and everything collapses to a point
(\cref{prop:nec-mono}).  The purpose of this paper is to say
precisely when they are safe, and to prove it.

\paragraph{Approach.}
We work over $\CICm$ (\cref{sec:calculus}), a stratified fragment
of CIC with an impredicative universe of propositions, two
predicative universes, primitive connectives, and parametric
inductive data types --- rich enough to express polymorphic
premises, their instances, and all the guard phenomena at issue,
while excluding proof-dependent data (we prove proofs are
hermetically confined in the fragment,
\cref{lem:proof-confinement}).  The key methodological decision is
semantic: we fix a set-theoretic standard model of the environment
(\cref{sec:semantics}, following
\cite{Werner1997,Aczel1998,LeeWerner2011}) and prove a
\emph{bidirectional truth lemma} (\cref{thm:truth}) for a guarded
translation in which the guard predicate is interpreted as
\emph{set membership} (\cref{sec:translation}).  Every
transformation of the paper is then measured against one invariant:
the translated axioms must remain true in the induced first-order
structure, and truth of the translated goal must continue to imply
validity of the original goal (\cref{def:si}).  This
model-theoretic setup is exactly what the erasure analysis needs
and what the existing proof-theoretic soundness framework
\cite{Czajka2016} cannot provide; it also matches the actual use of
\emph{classical} ATPs, for which minimal-logic proof reconstruction
does not apply.  The guarantee it delivers is semantic ---
optimised ATP successes imply the goal is true --- and combines
with the kernel-checked reconstruction phase, which makes the
end-to-end system safe unconditionally
(\cref{rem:absolute-safety}).

\paragraph{Contributions.}
\begin{enumerate}
\item A self-contained, model-theoretic soundness theorem for a
  guarded CIC-to-FOL translation in the style of CoqHammer
  (\cref{thm:truth,thm:base-soundness}), including a polarity
  analysis showing which of the implemented translation's guard
  forms are validated by the semantics and which are not
  (\cref{rem:polarity}).
\item A formulation of heuristic monomorphisation as premise
  instantiation at the $\CICm$ level --- instantiate, normalise,
  retranslate --- with: soundness independent of the
  instance-selection heuristic (\cref{thm:mono-sound});
  conservativity of the keep-generic policy
  (\cref{prop:mono-conservative}); incompleteness of the
  drop-generic policy (\cref{prop:drop-incomplete}); coverage of
  higher-order instances (\cref{rem:tier2}); and a proof that no
  computable complete instance-selection strategy exists for
  $\CICm$ (\cref{thm:undecidable}), adapting the
  Bobot--Paskevich combinatory-logic reduction
  \cite{BobotPaskevich2011} to a semantic consequence relation,
  with a full decidability analysis of the instance fragment.
\item \emph{Signature specialisation} (\cref{sec:specialization}):
  a uniform per-occurrence renaming --- type-argument fusion, guard
  specialisation $\hastype(x, \tau) \mapsto \guard{\tau}(x)$, arity
  flattening, predicate flattening --- proved sound by one
  interpretation principle (\cref{thm:spec-sound}) and proved
  \emph{faithful} (provability-preserving in both directions,
  \cref{thm:spec-faithful}) on the first-order core, \emph{without
  the bridging axioms} the current implementation emits.
\item \emph{Monotonicity-based guard erasure}
  (\cref{sec:erasure}): a formulation of guard erasure for the
  fully monomorphised, fully specialised problems; the reduction of
  its soundness to \emph{inflatability}
  (\cref{thm:erasure-from-inflatable}) with unconditional
  proof-preservation (\cref{prop:erasure-complete}); a sorted-core
  extraction and relativisation theorem
  (\cref{thm:relativisation}); full adaptations of the cloning and
  merging model constructions
  (\cref{lem:cloning,thm:erasure-main}); an \emph{internal}
  infinity criterion for inductive types driven by the
  translation's own injectivity and discrimination axioms
  (\cref{lem:ind-infinity}); and proofs that both side conditions
  --- witnessed inhabitation and monotonicity --- are necessary
  (\cref{prop:nec-inhab,prop:nec-mono}), the latter subsuming the
  known counterexample of \cite[\S5.6]{CzajkaKaliszyk2018}
  (\cref{rem:jar-counterexample}).
\item The composed pipeline with its end-to-end soundness theorem
  and an account of exactly where completeness is traded away
  (\cref{sec:pipeline}).
\end{enumerate}

All proofs are given in full, except for standard metatheory of
Pure Type Systems and inductive types
(\cref{prop:metatheory}) and the existence of set-theoretic models
(\cref{prop:model-exists}), which are quoted with references.

\paragraph{Design constraint.}
Throughout, transformations are \emph{shallow}: per-occurrence
manipulations of formulas, without mediating axiomatic bridges
between old and new symbols.  This constraint, adopted from the
CoqHammer development guidelines, is not cosmetic: bridging axioms
reintroduce the generic apparatus into every problem and defeat the
purpose of specialisation.  The faithfulness and erasure theorems
show that on the fragment that matters --- fully monomorphised,
essentially first-order problems --- nothing is lost by shallowness.

\paragraph{Organisation.}
\Cref{sec:calculus} defines $\CICm$ and proves the two fragment
lemmas.  \Cref{sec:semantics} fixes the semantics.
\Cref{sec:translation} defines the translation and proves the truth
lemma and base soundness.  \Cref{sec:mono,sec:specialization,%
sec:erasure} develop the three optimisations, and
\cref{sec:pipeline} composes them.  \Cref{sec:related} discusses
related work.
